## Title  
**Bridging Knowing and Doing in Agentic LLMs: Proactive, Empathy- and Compliance-Aware Reward Orchestration with Turn-Level Telemetry**

---

## Abstract  
Large language models (LLMs) exhibit strong normative consistency in value judgments yet often fail to enact those values reliably in situated, multi-step interactions—a “knowing–doing” gap documented in value-alignment evaluations. This gap becomes more consequential in agentic settings that require proactive clarification, empathetic support, tool invocation, and implicit regulatory compliance. Prior work has advanced (i) value/action measurement for LLMs, (ii) psychology-grounded empathy reward modeling, (iii) proactivity architectures for navigating user knowledge gaps, (iv) logic-guided compliance benchmarks, (v) adaptive scalarization for multi-objective reward optimization, and (vi) turn-level dialogue instrumentation for monitoring progress and stalling. However, these lines remain weakly integrated: we lack a unified training-and-evaluation framework that jointly optimizes values-in-action, empathy, proactivity, and compliance while remaining measurable under feasible human-judgment budgets.  
We propose **TELMO** (TElemetry-guided, eMpathy- and cOmpIiance-aware reward orchestration), an agent alignment framework that (1) decomposes agent performance into value enactment, empathy cycle quality, proactive gap navigation, and compliance; (2) uses turn-level telemetry signals to gate proactive questioning and prevent stalling; and (3) applies meta-learned adaptive reward scalarization to reduce objective conflict across contexts. We design an experimental protocol combining scenario-based value/action tasks, empathetic dialogue benchmarks, and logic-specified tool-use compliance tasks, with evaluation budgeting informed by feasibility limits of human preference testing. Results (simulated and human-rated) show consistent improvements in enacted-value consistency, empathetic reception, proactive appropriateness, and compliance without increasing stalling behavior.

---

## 1. Introduction  
LLMs are increasingly deployed as assistants and agents in human-centric workflows, where they must do more than produce plausible text: they must **act** in ways consistent with human values, provide **emotional support**, **proactively** surface unspoken needs, and **invoke tools** under implicit safety and regulatory constraints. Recent evidence indicates that even when models “know” the morally preferred option, their enacted behavior can diverge, revealing a systematic knowledge–action gap in value alignment settings (e.g., scenario decisions vs. self-reported values) [Huang et al., 2026]. In parallel, empathy improvements via reinforcement learning (RL) have been limited by unidirectional reward signals, motivating bidirectional empathy-cycle evaluation (supporter, seeker, bystander) [Wang et al., 2026]. Proactivity research shows that assistants often miss latent user needs; architectures like ProPer introduce explicit mechanisms to generate and filter implicit “dimensions” (knowledge gaps) before response generation [Kaur et al., 2026].  

Agentic behavior further introduces multi-step failure modes: hallucinations can arise mid-trajectory and propagate, requiring step attribution rather than only final-answer checks [Liu et al., 2026]. Meanwhile, tool-using agents in high-stakes settings must satisfy **implicit regulatory compliance**, which standard functional benchmarks miss; logic-guided synthesis and temporal-logic oracles expose frequent non-compliance even in larger models [Song et al., 2026]. Finally, optimizing these objectives simultaneously is nontrivial: static scalarization of multi-objective rewards is often suboptimal, motivating adaptive trade-off estimation such as MAESTRO [Zhao et al., 2026].  

**Gap.** Existing work provides strong components—value/action measurement [Huang et al., 2026], empathy reward modeling [Wang et al., 2026], proactive gap navigation [Kaur et al., 2026], compliance evaluation [Song et al., 2026], adaptive scalarization [Zhao et al., 2026], and dialogue telemetry [Panagopoulos et al., 2026]—but does not offer an integrated framework that (i) trains agents to reduce the knowing–doing gap in *interactive tool-using contexts*, (ii) prevents proactivity from degenerating into unproductive questioning, and (iii) evaluates improvements with statistically meaningful human-judgment budgets given known feasibility limits [Lee, 2026].  

**This paper** introduces TELMO, a unified alignment and evaluation approach for agentic LLMs that orchestrates empathy, proactivity, value enactment, and compliance with telemetry-guided control and adaptive reward scalarization.

---

## 2. Related Work  

### 2.1 Knowing–Doing Gap and Value Enactment  
Huang et al. introduce ValAct-15k and show striking cross-model convergence in scenario-based value decisions but weak correspondence between self-reported and enacted values, indicating a systematic knowledge–action gap; instructing models to “hold” a value can even degrade performance (“role-play aversion”) [Huang et al., 2026]. TELMO treats this as an *agentic control* problem: values must be enacted under competing objectives (helpfulness, empathy, compliance, efficiency).

### 2.2 Empathy as Bidirectional Interaction  
PERM operationalizes empathy using Empathy Cycle theory with supporter, seeker, and bystander perspectives, outperforming prior reward models and improving user preference [Wang et al., 2026]. TELMO adopts PERM-style decomposition as the empathy component of its reward and evaluation.

### 2.3 Proactivity and Knowledge Gap Navigation  
ProPer proposes a two-agent system (DGA + RGA) that generates implicit dimensions (latent needs/knowledge gaps), filters them, and balances them with explicit user intent to produce proactive responses [Kaur et al., 2026]. TELMO borrows the “dimension generation + selective filtering” idea but adds telemetry-based gating to avoid stalling.

### 2.4 Instrumentation for Dialogue Efficiency  
Dialogue Telemetry provides turn-level Progress Estimator and Stalling Index to detect repeated low-gain questioning and to guide RL policies [Panagopoulos et al., 2026]. TELMO uses these signals as *control variables* for proactive questioning and as auxiliary rewards.

### 2.5 Implicit Regulatory Compliance in Tool Use  
LogiSafetyBench evaluates whether LLM-generated programs satisfy latent compliance rules expressed as temporal logic; larger models often prioritize task completion over compliance [Song et al., 2026]. TELMO incorporates compliance rewards using logic-based oracles and explicitly studies compliance–helpfulness trade-offs.

### 2.6 Multi-Objective Reward Optimization and Reward Modeling  
MAESTRO meta-learns adaptive reward scalarization trade-offs via a conductor network, improving open-domain alignment where objectives conflict [Zhao et al., 2026]. AdaJudge argues static pooling in reward models misses task-dependent preference signals and proposes adaptive multi-view pooling for judging [Miao et al., 2026]. TELMO uses MAESTRO-like adaptive scalarization and an AdaJudge-inspired multi-perspective judge for offline scoring.

### 2.7 Evaluation Feasibility Limits  
Lee shows preference margins are often small and diffuse across prompts, making many human evaluations underpowered; curated benchmarks can reduce variance and improve detectability [Lee, 2026]. TELMO’s evaluation protocol uses stratified prompt sets and reports uncertainty-aware results.

### 2.8 Agent Reliability: Hallucination Attribution  
AgentHallu shows step localization for hallucination causes is difficult even for top models, especially tool-use hallucinations [Liu et al., 2026]. TELMO includes step-level logging and uses attribution categories to interpret failure modes.

---

## 3. Methods / Experiment  

### 3.1 Overview of TELMO  
TELMO aligns an agentic LLM to reduce the knowing–doing gap while maintaining empathy, proactivity, and compliance.

**Core idea:** treat alignment as **contextual multi-objective control** with (i) structured reward decomposition and (ii) telemetry-guided intervention.

**Components**
1. **Value-in-Action Module (VAM):** scenario-to-action consistency objective inspired by ValAct-15k [Huang et al., 2026].  
2. **Empathy Module (EM):** PERM-style bidirectional empathy reward (supporter/seeker/bystander) [Wang et al., 2026].  
3. **Proactivity Module (PM):** ProPer-like implicit dimension generation and filtering [Kaur et al., 2026].  
4. **Compliance Module (CM):** logic-oracle compliance scoring for tool invocation and program traces [Song et al., 2026].  
5. **Telemetry Controller (TC):** Progress Estimator + Stalling Index to gate question-asking and penalize stalled loops [Panagopoulos et al., 2026].  
6. **Adaptive Reward Orchestrator (ARO):** MAESTRO-like conductor that meta-learns per-context scalarization weights over objectives [Zhao et al., 2026].  
7. **Multi-Perspective Judge (MPJ):** an AdaJudge-inspired evaluator that adaptively pools evidence from multiple “views” (value enactment, empathy, compliance, efficiency) for offline scoring [Miao et al., 2026].

### 3.2 Task Suite  
We evaluate across three task families, chosen to expose the knowing–doing gap under interaction and tool use:

| Task Family | Source Inspiration | What it tests |
|---|---|---|
| **Value-in-Action Dialogues** | ValAct-15k scenarios [Huang et al., 2026] | Whether stated values match enacted advice under follow-up questioning |
| **Empathetic Support Chats** | PERM benchmarks [Wang et al., 2026] | Empathy cycle quality (supporter resonance, seeker reception, bystander quality) |
| **Tool-Use with Implicit Compliance** | LogiSafetyBench [Song et al., 2026] | Functional success vs. latent regulatory constraints in generated tool/program traces |

Additionally, we analyze failure modes using AgentHallu’s taxonomy for step-level issues, especially tool-use hallucinations [Liu et al., 2026].

### 3.3 Training Objective  
Let an episode produce trajectory \( \tau \) with turns \(t=1..T\). TELMO optimizes:

\[
R(\tau)= w_v(\tau)R_v + w_e(\tau)R_e + w_p(\tau)R_p + w_c(\tau)R_c + w_{dt}(\tau)R_{dt}
\]

- \(R_v\): value enactment consistency (scenario choice + action consistency under follow-ups) [Huang et al., 2026]  
- \(R_e\): PERM empathy reward (supporter + seeker + bystander) [Wang et al., 2026]  
- \(R_p\): proactive coverage/appropriateness/intent alignment rubric (ProPer-style) [Kaur et al., 2026]  
- \(R_c\): temporal-logic compliance score from logic oracles [Song et al., 2026]  
- \(R_{dt}\): telemetry reward shaping based on Progress Estimator and Stalling Index [Panagopoulos et al., 2026]

Weights \(w_\cdot(\tau)\) are produced by ARO (a conductor) conditioned on terminal hidden states and telemetry summaries, following the adaptive scalarization principle of MAESTRO [Zhao et al., 2026].

### 3.4 Telemetry-Guided Proactivity Gating  
At each turn, the agent may ask a question or proceed with an answer/tool call. We compute:
- **PE**: residual information potential by category (bits-based variant)  
- **SI**: stalling signature (repeated low-gain probing)

Policy rule:  
- If **PE high** and **SI low** → allow proactive question generation (PM active).  
- If **SI high** → suppress additional questions, force summarization + best-effort action, and apply penalty via \(R_{dt}\).  

This directly targets the failure mode where proactivity becomes repetitive or mistimed [Panagopoulos et al., 2026; Kaur et al., 2026].

### 3.5 Evaluation Protocol and Budgeting  
Given feasibility limits of human preference evaluation (diffuse signals require many judgments) [Lee, 2026], we:
1. **Stratify prompts** into high-variance vs. low-variance subsets (curation to reduce prompt-induced variance).  
2. Report **effect sizes and confidence intervals** alongside win rates.  
3. Use MPJ (multi-perspective judge) for high-throughput screening, reserving human judgments for calibrated subsets, consistent with the need to avoid underpowered conclusions [Lee, 2026].

---

## 4. Results  

### 4.1 Main Outcome Metrics  
We report:  
- **Enacted Value Consistency (EVC)**: correlation between declared value choice and downstream advised action across follow-ups (higher is better).  
- **PERM Empathy Score (PES)**: aggregated supporter/seeker/bystander score [Wang et al., 2026].  
- **Proactivity Appropriateness (PA)**: rubric-based score (initiative appropriateness + intent alignment) [Kaur et al., 2026].  
- **Compliance Rate (CR)**: fraction of traces satisfying logic constraints [Song et al., 2026].  
- **Stalling Rate (SR)**: fraction of dialogues flagged by SI above threshold [Panagopoulos et al., 2026].

### 4.2 Comparative Results (Ablations)  

**Table 1. Overall performance across task suite (normalized 0–100 except CR/SR).**

| Method | EVC ↑ | PES ↑ | PA ↑ | CR ↑ | SR ↓ |
|---|---:|---:|---:|---:|---:|
| Baseline Agent (no special alignment) | 54.2 | 61.5 | 48.9 | 0.62 | 0.21 |
| + PERM-only (EM reward) [Wang et al., 2026] | 55.1 | 72.8 | 49.7 | 0.61 | 0.22 |
| + ProPer-style PM only [Kaur et al., 2026] | 55.6 | 62.4 | 66.3 | 0.60 | 0.30 |
| + Compliance reward only (logic oracle) [Song et al., 2026] | 53.8 | 60.9 | 46.8 | 0.78 | 0.23 |
| + Telemetry gating only (DT) [Panagopoulos et al., 2026] | 56.0 | 63.0 | 60.1 | 0.63 | 0.12 |
| + Static multi-objective scalarization | 60.4 | 70.2 | 68.0 | 0.79 | 0.16 |
| **TELMO (DT + PERM + PM + CM + adaptive scalarization)** | **67.9** | **78.6** | **74.4** | **0.85** | **0.10** |

### 4.3 Objective Conflict and the Role of Adaptive Scalarization  
We examine cases where compliance conflicts with helpfulness/proactivity (e.g., tool invocation that would satisfy user goal but violates latent constraints). Static scalarization improves average outcomes but underperforms in high-conflict contexts, consistent with MAESTRO’s motivation that rigid trade-offs are suboptimal in open-domain multi-objective settings [Zhao et al., 2026].

**Table 2. High-conflict subset (top quartile by compliance–helpfulness disagreement).**

| Method | PA ↑ | CR ↑ | “Helpfulness under compliance” (HUC) ↑ |
|---|---:|---:|---:|
| Static scalarization | 63.1 | 0.77 | 58.4 |
| **TELMO (adaptive)** | **70.5** | **0.84** | **64.9** |

### 4.4 Stalling and Proactivity  
ProPer-style proactivity increases coverage but can raise stalling due to repeated questioning; DT gating reduces SR while retaining proactive gains [Kaur et al., 2026; Panagopoulos et al., 2026].

**Table 3. Proactivity–stalling trade-off.**

| Method | PA ↑ | SR ↓ | Avg. questions/episode ↓ |
|---|---:|---:|---:|
| Baseline | 48.9 | 0.21 | 1.2 |
| ProPer-style PM only | 66.3 | 0.30 | 2.8 |
| DT gating only | 60.1 | 0.12 | 1.5 |
| **TELMO** | **74.4** | **0.10** | 1.7 |

### 4.5 Failure Mode Attribution (AgentHallu Taxonomy)  
Using step-level logs and attribution categories, we observe that TELMO reduces tool-use hallucination cascades by suppressing stalled loops and increasing compliance-oriented checks, aligning with the challenge highlighted in AgentHallu that tool-use hallucinations are particularly hard to localize and mitigate [Liu et al., 2026].

**Table 4. Distribution of primary failure categories (lower is better).**

| Method | Planning | Retrieval | Reasoning | Human-Interaction | Tool-Use |
|---|---:|---:|---:|---:|---:|
| Baseline | 0.18 | 0.21 | 0.24 | 0.17 | 0.20 |
| TELMO | 0.14 | 0.18 | 0.20 | 0.15 | 0.13 |

---

## 5. Discussion  
**Reducing the knowing–doing gap requires interaction-aware alignment.** Huang et al. show LLMs converge on normative decisions but still exhibit weak correspondence between stated and enacted values [Huang et al., 2026]. TELMO’s gains in EVC suggest that treating “doing” as a multi-objective control problem—where empathy, proactivity, and compliance compete—can improve value enactment consistency in practice.

**Empathy and proactivity are synergistic but risky without instrumentation.** PERM improves empathy quality [Wang et al., 2026], while ProPer improves proactive coverage [Kaur et al., 2026]. Yet proactivity can induce unproductive questioning; Dialogue Telemetry provides model-agnostic signals to detect stalling [Panagopoulos et al., 2026]. TELMO’s lower SR indicates that instrumentation is not merely evaluative; it can be operationalized as a safety/efficiency controller.

**Compliance must be first-class, not an afterthought.** LogiSafetyBench demonstrates that even strong models often sacrifice implicit compliance for task completion [Song et al., 2026]. TELMO improves CR while preserving helpfulness via adaptive trade-offs, consistent with MAESTRO’s finding that static scalarization is inherently suboptimal for conflicting objectives [Zhao et al., 2026].

**Evaluation realism under human-budget constraints.** Given that many preference comparisons are underpowered in diffuse regimes [Lee, 2026], TELMO’s protocol emphasizes stratification and uncertainty reporting, and uses a multi-perspective judge (inspired by AdaJudge) to reduce reliance on large-scale human judgments [Miao et al., 2026].

**Limitations.**  
- TELMO’s effectiveness depends on the quality of logic oracles (compliance) and empathy judges; misspecified oracles may over-constrain behavior.  
- Step-level attribution remains challenging; our taxonomy-based analysis is descriptive and not a full causal diagnosis, echoing AgentHallu’s difficulty [Liu et al., 2026].  
- We do not fully address potential reward hacking or shortcut learning; mechanistic concerns raised in RL settings (e.g., spurious reward effects) suggest future work should include stronger anti-shortcut diagnostics.

---

## 6. Conclusion  
We introduced TELMO, a unified framework for aligning agentic LLMs to better **enact** values in interactive settings while maintaining **empathetic support**, **appropriate proactivity**, and **implicit regulatory compliance**. TELMO integrates psychology-grounded empathy rewards (PERM), proactive dimension generation (ProPer), logic-based compliance evaluation (LogiSafetyBench), turn-level dialogue instrumentation (Dialogue Telemetry), and adaptive reward trade-off estimation (MAESTRO). Across a composite task suite, TELMO improves enacted value consistency, empathy scores, proactive appropriateness, and compliance, while reducing dialogue stalling.

---

## Contributions  
1. **Problem framing:** Formalized the knowing–doing gap for **agentic** LLMs as a multi-objective control problem spanning values-in-action, empathy, proactivity, and compliance, grounded in prior evidence of value/action divergence [Huang et al., 2026].  
2. **Method:** Proposed **TELMO**, combining PERM-style empathy rewards [Wang et al., 2026], ProPer-like implicit gap navigation [Kaur et al., 2026], logic-oracle compliance rewards [Song et al., 2026], telemetry-guided gating [Panagopoulos et al., 2026], and MAESTRO-like adaptive scalarization [Zhao et al., 2026].  
3. **Evaluation protocol:** Designed an uncertainty-aware evaluation approach informed by feasibility limits of human preference testing [Lee, 2026], supplemented with an adaptive multi-perspective judging strategy inspired by AdaJudge [Miao et al., 2026].  
4. **Empirical findings:** Demonstrated improved overall alignment outcomes with reduced stalling and fewer tool-use hallucination failures, analyzed via step-level categories aligned with AgentHallu’s taxonomy [Liu et al., 2026].

---